
\section{Процедура и дополнительные результаты CEM}\label{app:cem}

\subsection{Обучение и линейные пробы}

Использован код CEM на основе коммита 26f027dd6f0c. Параметры повторения: предварительно обученная ResNet-34, SGD с learning rate 0.01, momentum 0.9 и weight decay 0.00004; batch size 128; до 300 эпох, early stopping по val\_loss с patience 15; вес потерь понятий 5; размерность вектора одного понятия 16. Пять повторов различаются случайным началом и выбором 28 из 112 атрибутов. Деления фотографий на train, validation и test берутся из обработанного набора авторов.

Для CEM с понятиями при обучении используется замена предсказанной вероятности целевой меткой текущего условия с вероятностью 0.25. В условии с настоящей разметкой это метка понятия, а в условии со случайным кодом — соответствующий бит кода. В CEM без обучения на понятиях вес потерь по A равен нулю и подмена выключена. Поэтому сравнение с этим базовым уровнем меняет и режим подмены, а не только один коэффициент потерь.

Основная проба — логистическая регрессия отдельно для каждого B. Координаты стандартизуются по средним и стандартным отклонениям train; регуляризация выбирается по точности на validation. Устанавливается max\_iter = 5000. Дополнительно используется проба на исходных координатах с C = 1.0. При единственном классе в train применяется постоянный ответ. Дополнительная метрика average precision (AP) оценивает, насколько высоко модель ставит положительные ответы при ранжировании; чем она выше, тем лучше. AP усредняется по признакам, для которых в test есть как минимум пять положительных наблюдений; правило важно при сопоставлении метрик. Сетка C и все конкретные настройки сохранены в коде пробы, указанном в реестре.

В каждом сравнении сначала вычисляется разность метрик внутри повтора, затем среднее этих пяти разностей. 95\%-й интервал: средняя разность $\pm$ 2.776 $\times$ стандартное отклонение разностей / $\sqrt{}$5. Аналогично вычисляются доверительные интервалы средних значений метрик для отдельных условий обучения. Пять повторов не являются пятью независимыми выборками видов или фотографий. Фотографии и виды в основном эксперименте одни и те же; интервалы условны на выбранном наборе и процедуре.

{\small\begin{longtable}{>{\raggedright\arraybackslash}p{0.303\linewidth}>{\raggedright\arraybackslash}p{0.303\linewidth}>{\raggedright\arraybackslash}p{0.303\linewidth}}\toprule
Модель и обучение & Средняя точность B, \% & 95\%-й интервал средней точности\\\midrule\endhead
CEM, только виды & 92.715 & [92.364; 93.066]\\
Hybrid-CBM 448, настоящие A & 94.763 & [94.147; 95.378]\\
CEM, настоящие A & 95.123 & [94.456; 95.790]\\
CEM, случайный код, согласованные веса & 95.392 & [95.218; 95.567]\\
\bottomrule\end{longtable}}

\textbf{Таблица A1.} Основная настроенная проба. Интервал разности вычисляется парно и не получается вычитанием концов интервалов средних точностей двух условий.

\subsection{Hybrid-CBM и настройки пробы}

Описание статьи задаёт размерность Hybrid-CBM 448: 28 вероятностей и 420 свободных координат с Leaky-ReLU. Конфигурация из исходного кода использует 476 координат (28 + 448) без этой нелинейности. Обе реализации исследованы; основная статья использует вариант 448.

{\small\begin{longtable}{>{\raggedright\arraybackslash}p{0.221\linewidth}>{\raggedright\arraybackslash}p{0.221\linewidth}>{\raggedright\arraybackslash}p{0.221\linewidth}>{\raggedright\arraybackslash}p{0.221\linewidth}}\toprule
Разность & Исходные координаты, C = 1, п. п. & Настроенная проба, п. п. & AP настроенной пробы\\\midrule\endhead
CEM - Hybrid-CBM 448 & +0.26 [-0.31; 0.83] & +0.36 [-0.18; 0.90] & +0.015 [0.002; 0.029]\\
CEM - Hybrid-CBM 476 & +0.87 [-0.16; 1.90] & +1.11 [0.15; 2.06] & +0.029 [-0.001; 0.058]\\
\bottomrule\end{longtable}}

\textbf{Таблица A2.} 95\%-е доверительные интервалы средних разностей по пяти парным повторам. Авторский код проб не найден; первая колонка не объявляется точным воспроизведением неизвестной реализации авторов.

\subsection{Случайный код и веса потерь}

В случайном коде каждый бит разыгрывается независимо с вероятностью единицы 0.5, повторяющиеся коды видов переразыгрываются. Зерно генерации — 5000 плюс индекс повтора. Назначение не использует настоящие атрибуты и признаки изображений.

В исходной процедуре положительные ответы взвешиваются с учётом частот: вес равен отношению числа учебных фотографий к числу положительных ответов минус один. Для настоящих A средний вес по атрибутам меняется от 5.4 до 8.5 между повторами, а для случайного кода близок к единице. В нашем первоначальном варианте со случайным кодом веса пересчитывались по контрольным меткам, поэтому одновременно менялись разметка и веса в функции потерь. Дополнительный вариант использует веса настоящих A; именно он вынесен в основной текст. Одинаковые коэффициенты не гарантируют одинаковую сложность задач или величины градиентов.

{\small\begin{longtable}{>{\raggedright\arraybackslash}p{0.221\linewidth}>{\raggedright\arraybackslash}p{0.221\linewidth}>{\raggedright\arraybackslash}p{0.221\linewidth}>{\raggedright\arraybackslash}p{0.221\linewidth}}\toprule
Сравнение случайного кода с настоящими A & Точность, настроенная проба, п. п. & Точность, C = 1, п. п. & AP\\\midrule\endhead
Веса вычислены по случайному коду & +0.65 [0.23; 1.07] & +0.82 [0.31; 1.33] & +0.011 [0.005; 0.017]\\
Веса взяты из настоящих A & +0.27 [-0.27; 0.80] & +0.43 [-0.19; 1.05] & +0.005 [-0.005; 0.014]\\
\bottomrule\end{longtable}}

\textbf{Таблица A3.} Эффект согласования весов. При настроенной пробе согласование снижает точность проб для B на 0.38 [0.22; 0.55] процентного пункта по сравнению с вариантом, в котором веса пересчитываются по случайному коду.

План случайного кода записан 10 ноября, поправка о весах — 11 ноября до результатов соответствующего запуска, но после первоначальных результатов случайного кода. Правило предусматривало преимущество при положительной нижней границе в обеих пробах, эквивалентность при расположении обоих интервалов внутри $\pm$0.30 пункта и иначе «не решено». Граница $\pm$0.30 пункта была рабочим критерием плана, а не установленным порогом практической важности. Для согласованного варианта итог — «не решено». Это не отменяет наблюдения высокой абсолютной читаемости и отдельного сравнения с Hybrid-CBM, но ограничивает утверждение о разности с настоящими A.

\subsection{Свойства структурированных контролей}

{\small\begin{longtable}{>{\raggedright\arraybackslash}p{0.171\linewidth}>{\raggedright\arraybackslash}p{0.171\linewidth}>{\raggedright\arraybackslash}p{0.171\linewidth}>{\raggedright\arraybackslash}p{0.171\linewidth}>{\raggedright\arraybackslash}p{0.171\linewidth}}\toprule
Свойство & Настоящие A & Перетасовка & Чужой вид & По внешности\\\midrule\endhead
Различных наборов на 200 видов & 153–185 & 198–200 & 153–185 & 172–192\\
Средний модуль коэффициента $\varphi$ по парам атрибутов & 0.13–0.17 & 0.05–0.06 & 0.13–0.17 & 0.12–0.16\\
Предсказуемость A по исходным визуальным признакам, AP & 0.63–0.67 & 0.18–0.23 & 0.18–0.24 & 0.68–0.74\\
Предсказание B по самим меткам A, AP & 0.59–0.60 & 0.19–0.21 & 0.19–0.22 & 0.37–0.41\\
\bottomrule\end{longtable}}

\textbf{Таблица A4.} Диапазоны по пяти повторам; свойства меток оцениваются на уровне видов. Ориентир для AP B около 0.20. Таблица описывает различия контролей, но не выделяет независимые причинные эффекты этих свойств. Случайный код в эту таблицу не включён, поскольку эти оценки для него не сведены в сопоставимую таблицу. Из его построения известны 200 различных кодов и вероятность единицы 0.5 в каждой позиции; остальные характеристики здесь не оцениваются.

\subsection{Другие постановки CUB}

В опыте переноса на виды, не встречавшиеся при обучении сети, разность настоящих A и «чужого вида» составила +0.017 [0.004; 0.030] AP, а настоящих A и обучения без понятий — -0.002 [-0.008; 0.003]. Проверка модели, обученной сразу на всех 112 атрибутах, также не дала убедительного преимущества над базовым уровнем: +0.0038 [-0.0046; 0.0123]. Поэтому по этому опыту нельзя заключить, что в представлении отсутствует информация о B для новых видов: чувствительность выбранной процедуры не подтверждена.

В отдельном анализе проб на десяти делениях видов сравнивались внутренний вектор сети, обозначенный здесь F, и его объединение с настоящими ответами A: [F, A]. При добавлении одних и тех же настоящих ответов A ко входу пробы разность AP между сетями, обученными на настоящих метках и на контроле «по внешности», составила -0.0065 [-0.0122; -0.0008], но двухэтапная проба (STK в коде анализа) дала -0.0036 [-0.0082; 0.0009]. В этой двухэтапной пробе сначала предсказывают B по одним A, затем логит этого предсказания стандартизуют, умножают на коэффициент 30 и присоединяют к F как дополнительную координату; для учебных примеров оценка получена без обучения на их видах. Смена модели пробы меняет вывод; эти результаты не используются как доказательство отсутствия информации «сверх A». Деления повторно используют одни виды и не считаются независимыми выборками популяции.

\section{Контрольные ответы B: проверка измерения}\label{app:probe}

Для каждого B независимо переставляем ответы между видами. Затем обучаем новые линейные пробы на тех же замороженных представлениях. Перестановка сохраняет число положительных видов, но при разных размерах классов может менять долю положительных фотографий. Использована одна перестановка для каждого из пяти повторов, с зерном 9200 плюс индекс повтора.

{\small\begin{longtable}{>{\raggedright\arraybackslash}p{0.221\linewidth}>{\raggedright\arraybackslash}p{0.221\linewidth}>{\raggedright\arraybackslash}p{0.221\linewidth}>{\raggedright\arraybackslash}p{0.221\linewidth}}\toprule
Представление & Настоящие B, \% & Контрольные B, \% & Разность, п. п.\\\midrule\endhead
CEM, только виды & 92.72 & 86.98 & 5.74\\
CEM, настоящие A & 95.12 & 90.13 & 4.99\\
CEM, случайный код с согласованными весами & 95.39 & 92.10 & 3.29\\
Hybrid-CBM 448, настоящие A & 94.76 & 89.51 & 5.25\\
\bottomrule\end{longtable}}

\textbf{Таблица B1.} Настроенные линейные пробы, пять повторов. Разность настоящей и контрольной читаемости называется избирательностью данного измерения. Базовый уровень 81.8\% относится к настоящим B и не переносится автоматически на переставленные ответы.

Избирательность у CEM на настоящих A выше, чем при согласованном случайном коде, на 1.70 [1.41; 1.99] пункта, но ниже, чем при обучении только на видах, на 0.75 [0.48; 1.02]. Порядок моделей зависит от того, рассматриваем ли мы точность на настоящих B или её разность с контрольной задачей. Избирательность не раскладывает точность на независимые «семантическую» и «несемантическую» составляющие. Анализ был записан до результатов без отдельного правила итогового вывода и рассматривается описательно; устойчивость к другим перестановкам не оценивалась.

\section{Протокол обзора литературы}\label{app:review}

\subsection{Отбор}

Протокол записан 3 ноября 2026 года, поиск выполнен 10 ноября. Период публикаций — 2018–2026, рецензируемые работы и препринты. Требуются эксперимент с обучающим сигналом A, оценка на других свойствах или задачах B и утверждение о связи результата с содержанием A. Допускаются осторожные формулировки; слабые или неявные приписывания отмечаются как пограничные.

Исключаются обзоры и позиционные статьи без экспериментов, только апостериорные пробы без обучения на A, оценка только самих A или основной задачи. Недоступный текст учитывается отдельно. Начальный план предусматривал Semantic Scholar, arXiv, OpenReview и при доступности Google Scholar. Фактически 11 запросов выполнены через OpenAlex, один — Semantic Scholar; Google Scholar не использован.

{\small\begin{longtable}{>{\raggedright\arraybackslash}p{0.1\linewidth}>{\raggedright\arraybackslash}p{0.5\linewidth}>{\raggedright\arraybackslash}p{0.16\linewidth}>{\raggedright\arraybackslash}p{0.12\linewidth}}\toprule
Код & Точная строка фактического запроса & Источник & Число\\\midrule\endhead
L1-1 & \texttt{\string"concept bottleneck\string" held-out concepts} & OpenAlex & 50\\
L1-2 & \texttt{\string"concept embedding\string" \string"unseen concepts\string"} & OpenAlex & 50\\
L1-3 & \texttt{\string"concept supervision\string" representation interpretability probe} & OpenAlex & 50\\
L1-4 & \texttt{\string"concept discovery\string" \string"sub-concepts\string"} & OpenAlex & 50\\
L1-5 & \texttt{\string"concept-based models\string" \string"leakage\string" held-out} & Semantic Scholar & 50\\
L1-6 & \texttt{\string"concept bottleneck\string" \string"zero-shot\string" new/unseen classes} & OpenAlex & 50\\
L1-7 & \texttt{\string"concept alignment\string" probe \string"remaining concepts\string"} & OpenAlex & 4\\
L2-1 & \texttt{\string"human similarity judgments\string" alignment representations transfer} & OpenAlex & 50\\
L2-2 & \texttt{\string"brain data\string" fine-tuning representations \string"human-like\string"} & OpenAlex & 50\\
L2-3 & \texttt{\string"intermediate task\string" fine-tuning probing linguistic} & OpenAlex & 50\\
L2-4 & \texttt{\string"cognitive alignment\string" \string"fine-tuning\string" behavior} & OpenAlex & 50\\
L2-5 & \texttt{\string"neural data\string" \string"auxiliary\string" \string"regularization\string" representations} & OpenAlex & 50\\
\bottomrule\end{longtable}}

\textbf{Таблица C1.} Первые 50 результатов по релевантности либо все результаты, если их меньше. В исходном плане L1-6 записан с окончанием \texttt{new classes / unseen classes}; в фактическом отчёте — \texttt{new/unseen classes}. В таблице приведены буквальные строки запросов; кавычки относятся только к отмеченным фразам.

Дополнительно выполнен один шаг по цитирующим работам для CBM (Koh и др., 2020), CEM, AligNet (Muttenthaler и др., 2025) и HiCEM. Для каждой исходной статьи отбирались до 100 цитирующих работ с наибольшим числом цитирований. Дополнительно включались цитирующие работы 2025–2026 годов, в названии или аннотации которых упоминалась оценка на свойствах или задачах, не использованных при обучении. Поиск по запросам дал 554 записи, по цитированиям — 337.

{\small\begin{longtable}{>{\raggedright\arraybackslash}p{0.468\linewidth}>{\raggedright\arraybackslash}p{0.468\linewidth}}\toprule
Этап & Записей\\\midrule\endhead
Всего до удаления дублей & 891\\
После первого удаления дублей & 802\\
Исключено по названию и аннотации после повторного просмотра & 633\\
Передано на чтение текста & 169\\
Дополнительные дубли & 2\\
Уникальных кандидатов & 167\\
Недоступный текст & 23\\
Исключено по критериям & 89\\
Включено & 55\\
\bottomrule\end{longtable}}

\textbf{Таблица C2.} Схема отбора. Дубли сопоставлялись по идентификатору arXiv и нормализованному названию.

\subsection{Извлечение и проверка}

Агенты языковых моделей разбирали работы партиями по 10–15. Для каждой работы записаны одинаковые поля: A, B, модель и данные, цитата утверждения, место в тексте, наличие и результат контроля, другие сравнения и объём чтения. Цитата ограничена 25 словами. Часть работ просмотрена по ключевым разделам; объём зафиксирован отдельно. При последующей сверке исправлен пропущенный контроль gLocal. Независимая двойная экспертная классификация не выполнялась, согласие разметчиков не измерялось.

Каталог 55 работ воспроизводит извлечённые сведения с цитатами и местами утверждений. Пометка «пограничный случай» взята из итоговой таблицы отбора, остальные поля — из таблицы извлечения. В исходном CSV 56 включённых строк: одна работа фигурирует под двумя названиями. Строка с отметкой «ДУБЛЬ при сводке» не включается повторно. Ранее вручную найденные статьи вне систематического отбора не добавляются в знаменатель.

Категория «есть» описывает наличие замены сигнала при том же формате ответов и функции потерь, а не сохранение всех статистических свойств. Категория «не найден» ограничена проведённым чтением. При интерпретации долей следует учитывать и пограничные случаи, и процедуру извлечения.

\section{Дополнительный опыт на GoEmotions}\label{app:goemotions}

Используем 12 000 коротких сообщений, ограничение 64 токена. Основная задача — отнести сообщение к одной из четырёх групп, заданных квартилями распределения длины в словах. Основа \texttt{multi-qa-distilbert-cos-v1}: нижние четыре слоя заморожены, верхние два дообучаются; добавлены отдельные выходные модули для основной задачи и понятий, BCE с коэффициентом 1. Выполнено пять инициализаций.

A включает anger, annoyance, approval, confusion, desire, disapproval, disgust, grief, nervousness, pride, relief, sadness, surprise. B включает admiration, amusement, caring, curiosity, disappointment, embarrassment, excitement, fear, gratitude, joy, love, optimism, realization, remorse. Деление фиксировано, это не пять разных делений эмоций.

Читаемость здесь измеряется средней AP по эмоциям B (macro-AP). Для каждой эмоции используется линейная регрессия с квадратичным штрафом на веса (ridge), решение которой вычисляется аналитически. В варианте с известными метками A их добавляют ко входу пробы отдельным блоком с подбираемым масштабом. «Обмен строками» передаёт сообщению полный набор меток другого сообщения; «гиперплоскости» и «кластеры» строят метки из исходных признаков с согласованием частот. Дополнительно исследован ошибочный контроль «случайные эмоции», сохранявший число настоящих меток у каждого сообщения.

{\small\begin{longtable}{>{\raggedright\arraybackslash}p{0.468\linewidth}>{\raggedright\arraybackslash}p{0.468\linewidth}}\toprule
Обучение & Средняя AP на B\\\midrule\endhead
Исходная модель & 0.376\\
Только задача длины & 0.321\\
Настоящие A & 0.473\\
Гиперплоскости & 0.353\\
Обмен строками & 0.319\\
Кластеры & 0.310\\
Случайные эмоции с сохранением числа меток сообщения & 0.428\\
\bottomrule\end{longtable}}

\textbf{Таблица D1.} Пять инициализаций; последнее условие содержит связь с B через структуру разметки и не устраняет зависимость контрольных меток от B.

{\small\begin{longtable}{>{\raggedright\arraybackslash}p{0.303\linewidth}>{\raggedright\arraybackslash}p{0.303\linewidth}>{\raggedright\arraybackslash}p{0.303\linewidth}}\toprule
Настоящие A минус контроль & Обычная проба, AP & Проба с известными A, AP\\\midrule\endhead
Гиперплоскости & +0.120 [0.091; 0.150] & +0.126 [0.096; 0.156]\\
Обмен строками & +0.153 [0.111; 0.195] & +0.150 [0.109; 0.192]\\
Кластеры & +0.163 [0.157; 0.169] & +0.166 [0.158; 0.173]\\
\bottomrule\end{longtable}}

\textbf{Таблица D2.} Парные 95\%-е t-интервалы по пяти инициализациям. Средние были известны до плана досчёта интервалов 10 ноября; этот анализ не объявляется заранее неизвестным результатом.

Контроль «обмен строками» почти не предсказывается по входу (AP 0.025) и не выучивается моделью (0.026); это слабое основание для изоляции роли содержания. Гиперплоскости предсказываются по входу с AP 0.483 против 0.209 у настоящих A. Поэтому основное положительное сравнение проводится с гиперплоскостями, но и оно не меняет только одно свойство меток. «Случайные эмоции» сохраняют число положительных меток A у каждого сообщения. В частности, сохраняется наличие хотя бы одной метки A, которое связано с B; их высокий результат нельзя приписывать произвольным меткам вообще. При обмене целых строк такого выигрыша над обучением только на задаче нет: -0.001 [-0.018; 0.016] AP.

\section{Дополнительная перепроверка DreamSim}\label{app:dreamsim}

\subsection{Постановка}

Рассматриваем перенос одной ветви DreamSim — DINO ViT-B/16, cls-признаки — после обучения на человеческих сравнениях NIGHTS. Использован код авторов, коммит db4d16c; LoRA rank 16, alpha 32, dropout 0.2; learning rate 0.0003, batch 32, восемь эпох, параметр margin функции потерь 0.05. Берутся тройки с не менее чем шестью единогласными голосами и исходным делением train/validation/test. Это ограниченное повторение одной ветви, не всего ансамбля статьи.

Сравниваются исходная модель, обучение на человеческих ответах и обучение на контроле. Контроль выбирает более близкое к эталону изображение по косинусному расстоянию после случайной гауссовой проекции исходных DINO-признаков из 768 в 64 измерения. Он не использует ответы людей, но наследует структуру исходного представления и не объявляется лишённым любой семантики. Пять инициализаций, зерно проекции 9100 плюс индекс запуска. Проверочные ответы NIGHTS остаются настоящими.

Внешняя оценка — корреляция Спирмена (SROCC) с человеческими оценками качества изображений TID2013 и KADID-10k. Запланированная оценка BAPPS не выполнена из-за недоступности данных; изменение записано до основного запуска. Источник весов DINO и режим TF32 также зафиксированы в журнале процедуры. Побайтовая идентичность весов альтернативного источника с исходной загрузкой DINO не проверялась.

\subsection{Результаты}

{\small\begin{longtable}{>{\raggedright\arraybackslash}p{0.221\linewidth}>{\raggedright\arraybackslash}p{0.221\linewidth}>{\raggedright\arraybackslash}p{0.221\linewidth}>{\raggedright\arraybackslash}p{0.221\linewidth}}\toprule
Набор & Исходная модель & Человеческие ответы & Случайная проекция\\\midrule\endhead
TID2013, SROCC & 0.793 & 0.832 & 0.813\\
KADID-10k, SROCC & 0.841 & 0.831 & 0.853\\
NIGHTS test, точность \% & 90.13 & 94.07 & 89.25\\
\bottomrule\end{longtable}}

\textbf{Таблица E1.} Средние по пяти инициализациям для обучаемых условий; исходная модель фиксирована. NIGHTS соответствует обучающей области и не заменяет внешнюю оценку.

{\small\begin{longtable}{>{\raggedright\arraybackslash}p{0.303\linewidth}>{\raggedright\arraybackslash}p{0.303\linewidth}>{\raggedright\arraybackslash}p{0.303\linewidth}}\toprule
Набор & Человеческие ответы минус случайная проекция & Случайная проекция минус исходная модель\\\midrule\endhead
TID2013, SROCC & +0.020 [0.015; 0.024] & +0.020 [0.016; 0.023]\\
KADID-10k, SROCC & -0.022 [-0.027; -0.018] & +0.012 [0.008; 0.015]\\
NIGHTS test, п. п. & Не используется как основной внешний контраст & -0.88 [-1.49; -0.26]\\
\bottomrule\end{longtable}}

\textbf{Таблица E2.} 95\%-е t-интервалы по пяти инициализациям; сравнение обученных условий парное. Интервалы описывают вариацию обучения при фиксированных наборах изображений, не вариацию выбора изображений или оценщиков.

На TID2013 человеческие ответы дают дополнительный прирост относительно выбранного контроля, но контроль также улучшает исходную модель. На KADID-10k настоящие ответы не улучшают исходную модель и уступают контролю. По правилу, добавленному в план 11 ноября, разные выводы на двух наборах дают общий итог «не решено». Поэтому нельзя обобщить преимущество человеческих ответов на оба внешних набора или использовать данный опыт как однозначное опровержение переноса DreamSim.

Перемешанные и противоположные человеческим ответы существенно ухудшали результат относительно исходной модели. Такие разрушительные контроли сами по себе не отделяют полезность содержания от необходимости обучаться согласованной с входом задаче. Контроль со случайной проекцией добавили в план до получения его результатов; полный набор решений сохранён в реестре.

\section{Дополнительные сведения о HiCEM}\label{app:hicem}

Использован код авторов на коммите 4b665fa, CUB и CLIP ViT-L/14. Для каждого из четырёх вариантов выполнено три запуска. Обозначения в журналах: R — исходная процедура; Rtrue — настоящие A с отбором фотографий по настоящим общим меткам; P — контроль «чужой вид» с тем же отбором; C1 — поиск непосредственно по CLIP с тем же отбором. Во всех вариантах последующая HiCEM обучается на настоящих общих метках и найденных метках подпонятий; в P контрольные метки используются только при обучении исходной CEM. Общие понятия, внутри которых проводится поиск подробных признаков, далее называются родительскими. Для перестановочной проверки метки каждого подпонятия переставляются между видами внутри его родительского понятия. Число положительных видов сохраняется, а выбор лучшего кандидата повторяется на каждой перестановке. Сопоставление и перестановочная проверка используют разметку учебных фотографий на этапе поиска подпонятий. Одновременно проверяются 188 понятий. Для каждого выполняется K = 1000 рандомизаций; поправка Бенджамини—Хохберга (BH) применяется с номинальным уровнем q = 0.1. Для этой проверки берётся лучший кандидат каждого понятия банка до применения авторского порога 0.7 и разрешения конкуренции за кандидата. При подсчёте не требуется, чтобы значимое соответствие одновременно удовлетворяло авторскому правилу сопоставления; это два разных критерия. Рандомизация предполагает обменяемость видов внутри родительского понятия и меняет связи между подпонятиями. Гарантия контроля доли ложных открытий при произвольной зависимости этих проверок не доказана. Поэтому значимость соответствия определяется выбранной схемой рандомизации и правилом BH; альтернативные способы рандомизации не проверены.

Двухуровневый бутстрэп использует 2000 выборок: сначала повторно выбираются запуски, затем родительские понятия. При бутстрэпе повторно выбираются уже вычисленные индикаторы сопоставления и значимости; перестановочные p-значения и отбор BH заново не рассчитываются. Сравнение средней тестовой AUC последующей HiCEM между P и Rtrue использует отдельный бутстрэп по трём запускам. Эта AUC взята из результатов предсказания подпонятий на тестовых фотографиях и отличается от AUC сопоставления кандидатов на этапе поиска. Эти интервалы условны на данных, банке понятий и выбранной процедуре; при трёх запусках оценка хвостов распределения ограничена. Основной текст использует результаты 119/120/122 для R, 124/122/114 для Rtrue, 81/78/83 для P и 75/79/77 для C1.

Отдельный контроль заменяет банк подпонятий случайными группами видов. Средняя тестовая AUC последующей HiCEM, рассчитанная по процедуре авторов, находится в диапазоне 0.799–0.814; в раннем анализе он не дал значимых соответствий после рандомизационной проверки. Этот результат относится к используемой схеме рандомизации и не доказывает отсутствия любых понятий в представлении. Данные не позволяют распространять результат CUB на остальные наборы статьи HiCEM.


В дополнительном сравнении оставлены 51 подпонятие, найденное во всех запусках всех четырёх условий. Для каждого подпонятия, запуска и условия из наблюдаемой AUC сопоставления вычитается средняя AUC при перестановках его меток. Здесь используется AUC этапа поиска на учебных фотографиях, а не тестовая AUC последующей HiCEM. Разность этого превышения между Rtrue и P составляет +0.006 [-0.009; +0.018]. Такое сопоставление фиксирует набор понятий, но исключает находки, присутствующие лишь в части условий.
